\label{sec:cohorts}

\subsection{Cohorts}

\textbf{NADT-Prostate} (TCIA, DOI:10.7937/TCIA.JHQD-FR46): 39 patients, H\&E and ERG-stained
serial sections plus a PIN-4 cocktail, with real per-core Gleason score and phenotype
(benign/tumor) labels. Used to fit markers 1--3.
\textbf{PANDA} (Kaggle): 10{,}616 biopsy slides from two institutions (Karolinska, Radboud)
with official ISUP grade 0--5 labels; a stratified 1{,}200-slide subsample (1{,}137 after
tissue filtering) is used for external validation of markers 1--2.
\textbf{TCGA-PRAD} (GDC + cBioPortal \texttt{prad\_tcga\_pub}): 300 H\&E slides / 273 patients
with real Gleason score, PTEN copy-number, SPOP mutation, AR-activity score, and ERG fusion
status; also used, via a from-scratch relabeling of the GDC clinical API
(Section~\ref{sec:recurrence-transfer}), for recurrence-style outcome labels.
\textbf{SICAPv2}: used for diagnostic sanity checks (zero-shot text-prompt vs.\ linear-probe
comparison) reported in Section~\ref{sec:core-results}.
\textbf{LEOPARD} (\texttt{leopard.grand-challenge.org}): 508 patients, single institution
(Radboudumc), real biochemical-recurrence event and follow-up-time labels
(\texttt{case\_id}/\texttt{event}/\texttt{follow\_up\_years}; no clinical covariates in the
public release). Obtained from the dataset's public S3 bucket; subject to a publication
embargo pending the official challenge and baseline papers, which will be re-confirmed with
organizers before submission (see the Ethics and data-use statement).
\textbf{PRECISE} (Zenodo, DOI:10.5281/zenodo.20721779): 25 patients, one core-needle biopsy
session each except one patient with three sessions (27 imaging sessions total, confirmed
against the dataset's own participant manifest and README, which also confirms these are
core-needle biopsies, not resections -- the manifest's age field is explicitly documented as
``age at biopsy''), paired H\&E and CKAPM+racemase IHC whole-slide images with 24{,}387
pixel-level expert annotations across seven classes (background, tumor, benign gland, artifact,
HGPIN, intraductal carcinoma, atypical intraductal proliferation, stroma) and, separately, real
per-image pathologist Gleason score.

\subsection{Frozen encoders}

CONCH (\texttt{hf\_hub:MahmoodLab/conch}, ViT-B/16, 512-dim, \texttt{proj\_contrast=False})
and Virchow (\texttt{hf\_hub:paige-ai/Virchow}, ViT-H/14, 2560-dim = concatenated CLS token
and mean of patch tokens, requiring \texttt{mlp\_layer=SwiGLUPacked, act\_layer=SiLU}) are used
as frozen encoders; no weights are fine-tuned anywhere in this study. Tiles are 448$\times$448
(CONCH) or 224$\times$224 (Virchow). Whole-slide images without a usable multi-resolution
pyramid (LEOPARD, PRECISE) are tiled via windowed \texttt{zarr} reads against the
baseline-resolution page rather than full in-memory decoding.

For the completed Gate A grid, the frozen NADT representation contains all 463 H\&E slides
for phenotype and the prespecified 334-slide evaluable subset for Gleason; these are distinct
analysis frames and are not interchanged. TCGA-PRAD contributes the fixed PTEN, SPOP, and AR
frames. Marker 7 uses a fixed 270-patient TCGA-PRAD target, while the retained LEOPARD source
membership varies from 498 to 502 patients across the already-frozen configurations; their
common intersection is 498. The completed common-498 R2 analysis restricted all 60 cells to
those same 498 source patients (85 events) and retained the fixed 270-patient TCGA-PRAD target
(57 events). This closes the source-membership imbalance check without treating the correlated
encoder/scale settings as independent cohorts.

Tile selection is \emph{not} uniform across cohorts; each cohort's actual method, with its
tissue-fraction threshold and per-slide tile count, is as follows. \textbf{NADT-Prostate and
PANDA} (markers 1--3 fitting and PANDA zero-shot transfer): random-uniform rejection sampling
-- a random $(y_0,x_0)$ tile origin is drawn repeatedly (up to 200 attempts per slide) at a
fixed pyramid level ($\approx 0.88\,\mu$m/px, chosen to match SICAPv2's native
$\approx 1\,\mu$m/px scale) and kept if the fraction of non-white pixels (mean grayscale
intensity $<205$) is $\geq 0.35$, until 16 tiles/slide are collected (PANDA deliberately reuses
NADT's exact tile count for a fair zero-shot comparison). \textbf{TCGA-PRAD} (markers 4--7,
ERG-fusion status, and all additional molecular candidates): the same random-uniform
rejection method and $0.35$ tissue-fraction threshold, but at a target $\approx 0.88\,\mu$m/px
($\approx 0.97\,\mu$m/px for the Virchow scale-matched re-embedding used in the marker-7 cross-
encoder check) derived \emph{dynamically per slide} from each file's AppMag tag rather than a
fixed pyramid-level index (Supplementary S3 records why this matters); 16 tiles/slide (64 for
the Virchow scale-matched re-embedding, matching LEOPARD Virchow's tile count for that specific
comparison). \textbf{LEOPARD}: non-overlapping grid sampling (not random) restricted to each
slide's own provided tissue-segmentation mask, kept if the mask-fraction within a tile is
$\geq 0.5$ (a real segmentation mask, not a non-white-pixel heuristic), randomly subsampled to
64 tiles/slide when the grid yields more candidates. \textbf{PRECISE}: non-overlapping grid
sampling against the dataset's pixel-level expert annotation mask, at a $150\,\mu$m physical
window (native $\approx 0.243\,\mu$m/px, resized to 448$\times$448) chosen specifically because
the standard $\approx 0.88\,\mu$m/px scale produced windows too coarse to find any qualifying
single-class region in this cohort's small annotated areas; a grid cell is kept if one
non-background class occupies $\geq 70\%$ of it, capped at 20 tiles per class per image (not
per slide overall, since the goal is balanced representation across tissue classes rather than
a fixed total). Per-slide embeddings are the mean of the tiles collected under each cohort's own
method above; per-patient values, where used, are the mean over a patient's slide-level
embeddings (or, for PRECISE, session-level embeddings, Section~\ref{sec:core-results}).
Gate A shows marker-specific, mixed associations between physical scale and observed transfer
across correlated configurations. It does not isolate scale from sampling or encoder, and it
does not establish scale as an independent or causal source of degradation
(Section~\ref{sec:failure-modes}). Common-498 restriction did not explain or reverse marker
7's observed directions, but it remains a correlated sensitivity analysis rather than a
causal scale experiment or independent replication.

The Gate A sensitivity grid uses 16, 32, and 64 tiles per slide under five sampling seeds.
CONCH is evaluated at native $0.88$ and shared $1.76\,\mu$m/px; Virchow at native $0.44$ and
shared $1.76\,\mu$m/px. These are frozen correlated configurations, not independently sampled
patient cohorts. The older $0.97\,\mu$m/px marker-7 diagnostic described above is retained as
historical context only and is not the source of the completed Gate A summaries.
